\documentclass[
	headinclude=false,
	footinclude=false,
        12pt,
	]{scrartcl}
    
\usepackage[margin=25mm]{geometry}

\usepackage[T1]{fontenc}% moderne Schriftkodierung
\usepackage[utf8]{inputenc}% Sonderzeichen vieler verschiedener Sprachen gleichzeitig (€, ° usw.)
\usepackage[english,main=ngerman]{babel}% Primär Deutsch und sekundär Englisch
\usepackage[autostyle]{csquotes}% Anführungszeichen mit \enquote{} passen zur Sprache von babel
\usepackage{lmodern}% Schriftfamilie latin modern basierend auf computer modern verträgt sich besser mit fontenc T1
\usepackage{microtype}% Verbessert den Textsatz durch z.B. geringfügige Buchstaben-Skalierungen

\renewcommand{\familydefault}{\sfdefault}
    
\usepackage{scrhack}% löst Probleme mit floats
\usepackage[pdftex]{graphicx}

\usepackage{enumitem}
\setlist[itemize]{noitemsep, topsep=6pt}


\begin{document}
    \pagenumbering{gobble}
    \begin{center}
        {\LARGE\textbf{\mbox{Hochschulmeisterschaften Volleyball WS 25/26}}}
    \end{center}
    \vspace{1.5em}
{\Large\textbf{Spielregeln}}
\large
\begin{itemize}
    \item Maßgebend sind die offiziellen FIVB-Regeln 2025-2028.
    \item Pro Spiel werden zwei Sätze gespielt. Ein Satz geht solange bis die erste Mannschaft 15 Punkte erreicht.
    \item Die Entscheidung des Schiedsgerichts ist zu respektieren.
    \item Pro Satz darf jeder Spieler nur einmal ein- oder(!) ausgewechselt werden darf. 
    \item Eine Frau muss die gesamte Rotation mitspielen, das heißt sowohl im Vorder- als auch Hinterfeld. 
\end{itemize}   
\vspace{1cm}
 {\Large\textbf{Schiedsgericht}}
\begin{itemize}
    \item Für ein Schiedsgericht müssen mindestens 2 Personen anwesend sein. Davon übernimmt eine Person das 1. Schiedsgericht. Die zweite Person übernimmt das Punkteklappen und Schreiben.
    \item Wenn mehr Personen vor Ort sind, sollen alle Linien besetzt werden. 
    \item Der Spielberichtsbogen muss direkt nach dem Spiel zur Turnierleitung gebracht werden.
    \item Stört ein Ball eines anderen Spiels einen Spielzug, wird sofort Doppelfehler gepfiffen.
\end{itemize}  
\vspace{1cm}
 {\Large\textbf{Fair Play}}
\begin{itemize}
    \item Über die Spielregeln hinaus achten und respektieren wir jedes Teammitglied und jeden Gegner, indem wir unter anderem das jeweilige Leistungsniveau berücksichtigen. Dies gilt vor allem im Aufschlag und im Angriff.
    \item Fair Play bedeutet beispielsweise \ldots
    \begin{itemize}
        \item[\ldots] nicht um jeden Preis gewinnen zu wollen.
        \item[\ldots] eigene Teammitglieder aufzumuntern, wenn sie einen Fehler gemacht haben.
    \end{itemize}
    \item Bei groben Verstößen gegen diesen Fair Play-Gedanken kann im Ermessen der Turnierleitung in einem nächsten Spiel z.B. der Verlust des Aufschlagsrechts erfolgen oder auch ein Spielverbot.
\end{itemize}    


    
 \newpage
\begin{center}
        {\LARGE\textbf{\mbox{Hochschulmeisterschaften Volleyball WS 25/26}}}
    \end{center}
    \vspace{1.5em}
{\Large\textbf{Rules of the games}}
\large
\begin{itemize}
    \item The official FIVB rules 2025-2028 apply.
    \item Two sets are played per game. A set continues until the first team reaches 15 points.
    \item The decision of the referees must be respected.
    \item Each player may only be substituted in or out once per set. 
    \item A woman must play the entire rotation, i.e., both in the front and back court. 
\end{itemize}   
\vspace{1cm}
 {\Large\textbf{Schiedsgericht}}
\begin{itemize}
   \item  At least two people are required to fulfil refereeing duties. One of them will act as the first referee. The second person is responsible for keeping score and writing reports.
   \item If there are more people on site, all lines should be occupied. 
   \item The match report form must be given to the tournament management immediately after the match.
    \item If a ball from another match interferes with play, a double fault is immediately called.
\end{itemize}  
\vspace{1cm}
 {\Large\textbf{Fair Play}}
\begin{itemize}
    \item Beyond the rules of the game, we respect every team member and every opponent by, among other things, taking into account their respective performance levels. This applies especially to serving and attacking.
    \item Fair play means, for example \ldots
    \begin{itemize}
        \item[\ldots] not wanting to win at any cost.
        \item[\ldots] encouraging your own team members when they have made a mistake.
    \end{itemize}
    \item In the event of gross violations of this spirit of fair play, the tournament management may, at its discretion, impose penalties in the next game, such as loss of the right to serve or even a ban from playing.
\end{itemize}    
\end{document}
